\chapter{Roteiro e atividades para discussão}
Os blocos citados nesta parte correspondem às seções numeradas do script R. As respostas esperadas apoiam a condução docente e a discussão coletiva.

\textbf{27 de outubro de 2026 · Estatística Aplicada · RENAmazon/UEAP}

\section{Objetivo da aula}

Mostrar como uma pergunta ecológica se transforma em análise auditável: explorar dados, representar variabilidade e incerteza, formular hipóteses, examinar modelos e interpretar estimativas. O eixo conceitual é a comparação entre variação explicada e residual, que conecta ANOVA e regressão.

O script foi desenvolvido para o professor. As funções de infraestrutura podem ser carregadas previamente; durante a aula, priorize os dados, as fórmulas dos modelos, os gráficos e as interpretações. Reserve a implementação matricial HC3 e as rotinas leave-one-out para discussão avançada, se houver tempo.

\section{Percurso sugerido --- 180 minutos, adaptável}

\begin{table}[H]\centering\small
\begin{tabular}{p{0.28\linewidth}p{0.38\linewidth}p{0.24\linewidth}}
\toprule
Tempo & Condução & Seções e figuras\\
\midrule
0--15 min & Perguntas, desenho, unidade e auditoria & Seção 1\\
15--35 min & Exploração: antecipar diferenças sem calcular p & Seção 2; figuras 01 e 02\\
35--65 min & Construir SQ, QM e F; visualizar a hipótese nula & Seção 3; figuras 03 e 04\\
65--80 min & Confrontar pressupostos com dados e protocolo & Seção 4; figuras 05 e 06\\
80--105 min & Comparações de Tukey e magnitude & Seção 6; figuras 08 e 09\\
105--135 min & Regressão simples, F = t\textsuperscript{2} e dois intervalos & Seções 7--8; figuras 10--13\\
135--165 min & Regressão múltipla e teste F parcial & Seções 9--10; figuras 14--16\\
165--180 min & Redação do resultado e síntese & Seção 13\\
\bottomrule\end{tabular}\end{table}

\textbf{Extensões:} randomização (seção 5), interação/HC3 (parte da seção 10), sensibilidade/LOO (seção 11) e DBC (seção 12). Para uma aula de 120 minutos, concentre-se no DIC, Tukey e regressão simples, deixando regressão múltipla e extensões para outra sessão.

\section{Antes de iniciar}

\begin{itemize}
\item Abra o \texttt{.Rproj} e carregue a seção 0, que define as funções e o padrão gráfico.
\item Deixe \texttt{saidas/relatorio\_aula.html} aberto como apoio à projeção; todas as figuras já estão disponíveis.
\item Apresente o caráter simulado dos dados antes de mostrar resultados. As médias e efeitos não descrevem a Amazônia real.
\item Diferencie o DIC de decomposição, em que o tratamento foi alocado às unidades, do estudo de aves, em que os ambientes são categorias observadas.
\end{itemize}

\section{Atividade 1 --- Ler os dados como um pesquisador}

\textbf{Pergunta:} o que cada linha representa? Quais observações poderiam ser réplicas e quais seriam subamostras?

\textbf{Execução:} seção 1; examine \texttt{str()}, contagens por grupo e resumos.

\textbf{Resposta esperada:} 24 recipientes no DIC, oito por umidade; 36 localidades de aves, 12 por ambiente; 18 unidades no DBC. Os quatro pontos de escuta agregados não fornecem quatro réplicas independentes do ambiente. Se uma única câmara recebesse cada nível de umidade, o efeito de câmara ficaria confundido com tratamento.

\textbf{Ênfase oral:} independência é uma propriedade do desenho e do processo, não um resultado de Shapiro. O número de linhas não define sozinho o tamanho amostral efetivo.

\section{Atividade 2 --- Antecipar padrões com exploração gráfica}

\textbf{Pergunta:} qual grupo parece ter maior perda de massa? A dispersão é comparável? Os grupos são totalmente separados?

\textbf{Execução:} seção 2; figuras 01 e 02.

\textbf{Valores de referência:} médias de 21,683\%, 29,458\% e 39,306\% para Baixa, Média e Alta; DP de 4,402, 4,696 e 3,652 pontos percentuais.

\textbf{Resposta esperada:} gradiente de médias com sobreposição de observações individuais. A sobreposição não impede diferenças entre médias; a posição de um único recipiente não representa a média de um tratamento. Com oito observações por grupo, a forma do histograma é instável; pontos brutos são essenciais.

\textbf{Recurso gráfico a demonstrar:} alterar \texttt{binwidth}, a transparência e a largura do jitter. Manter a mesma escala permite comparar dispersão; mudar escalas pode esconder diferenças. O jitter altera apenas a posição horizontal de exibição, sem mudar os dados analisados.

\section{Atividade 3 --- Construir o teste F sem tratá-lo como caixa-preta}

\textbf{Hipóteses:} H0: \ensuremath{\mu}Baixa = \ensuremath{\mu}Média = \ensuremath{\mu}Alta. H1: pelo menos uma média difere. \ensuremath{\alpha} = 0,05, definido antes da análise.

\textbf{Execução:} seção 3; compare os segmentos da figura 03 e a distribuição da figura 04.

\textbf{Questões para a turma:}

\begin{enumerate}
\item Por que centralizar todas as observações na média geral produz mais variação que centralizá-las nas médias dos tratamentos?
\item Por que dividir as somas de quadrados por graus de liberdade antes de formar a razão?
\item Se H0 fosse verdadeira, o que esperar da razão QMtratamento/QMerro?
\item O F global informa qual par difere?
\end{enumerate}

\textbf{Referências numéricas:} F(2,21) = 34,181; F crítico \ensuremath{\approx} 3,467; p \ensuremath{\approx} 2,49 \ensuremath{\times} 10\textsuperscript{-}\textsuperscript{7}; eta\textsuperscript{2} \ensuremath{\approx} 0,765; omega\textsuperscript{2} \ensuremath{\approx} 0,734.

\textbf{Interpretação:} existe evidência contra a igualdade das três médias sob o modelo e o desenho. F alto expressa grande separação relativa à variabilidade residual, mas não substitui magnitude na escala da resposta. Eta\textsuperscript{2} não significa que esse tratamento explicará a mesma fração da variação em todos os ambientes reais.

\textbf{Frase a corrigir em aula:} “A chance de H0 ser verdadeira é 0,000000249.” O valor-p é calculado sob a nula e as demais suposições; não é probabilidade posterior da hipótese.

\section{Atividade 4 --- Diagnosticar antes de concluir}

\textbf{Execução:} seção 4; figuras 05 e 06.

\textbf{Perguntas:} o Q--Q tem desvios sistemáticos? Há dispersão residual claramente distinta? Há posição experimental associada ao resíduo? O que esses gráficos não conseguem testar?

\textbf{Condução:} explique os quatro painéis e os resultados de Shapiro e Fligner como informações complementares. Não use “p > 0,05” para certificar pressupostos. No DIC balanceado, as alavancagens são iguais; diferenças de Cook decorrem dos resíduos, e não de posições extremas de preditores contínuos.

\textbf{Sensibilidade:} Welch retorna F aproximado \ensuremath{\approx} 37,129 e p \ensuremath{\approx} 2,75 \ensuremath{\times} 10\textsuperscript{-}\textsuperscript{6}. As duas análises apontam evidência contra igualdade de médias, mas o professor não escolhe retrospectivamente o teste com menor p. Se a heterocedasticidade fosse relevante, Tukey clássico também precisaria ser revisto; trocar apenas o F não conserta toda a análise.

\section{Atividade 5 --- Ler Tukey como diferenças, não como letras}

\textbf{Execução:} seção 6; figuras 08 e 09.

\begin{table}[H]\centering\small
\begin{tabular}{p{0.28\linewidth}p{0.38\linewidth}p{0.24\linewidth}}
\toprule
Diferença & Estimativa, p.p. & IC95\% simultâneo\\
\midrule
Média - Baixa & 7,775 & 2,390 a 13,160\\
Alta - Baixa & 17,624 & 12,239 a 23,009\\
Alta - Média & 9,849 & 4,464 a 15,234\\
\bottomrule\end{tabular}\end{table}

\textbf{Perguntas:}

\begin{enumerate}
\item O que define o sinal de cada diferença?
\item O intervalo cobre zero? Qual é a magnitude compatível com os dados?
\item Por que o intervalo de uma média não substitui o intervalo da diferença?
\item A diferença entre 39,306\% e 21,683\% é de 17,624\% ou 17,624 pontos percentuais?
\end{enumerate}

\textbf{Resposta esperada:} todos os pares apresentam IC que excluem zero. O maior contraste é Alta - Baixa, com maior perda na umidade alta. A unidade é ponto percentual. Aumento relativo seria outro cálculo, com divisão por uma referência.

\textbf{Ênfase:} o procedimento controla a família dos três pares. Não apresentar uma não rejeição como equivalência. Não usar significância como medida de importância ecológica. Perda de massa é um indicador de decomposição; não mede, sozinha, toda a ciclagem de nutrientes.

\textbf{Recurso gráfico:} compare o forest plot de diferenças com os pontos brutos e as médias. Evite letras na primeira demonstração para manter a atenção em estimativa, incerteza e unidade.

\section{Atividade 6 --- Conectar ANOVA e regressão}

\textbf{Execução:} seções 7 e 8; figuras 10--13.

\textbf{Hipótese:} H0: \ensuremath{\beta}topografia = 0. A alternativa bilateral admite inclinação diferente de zero.

\textbf{Referências:} \ensuremath{\widehat H} = 1,77292 + 0,02601 \ensuremath{\times} topografia em metros; F(1,34) \ensuremath{\approx} 20,104; R\textsuperscript{2} \ensuremath{\approx} 0,372. O F é numericamente igual ao quadrado do t da inclinação, porque há um único coeficiente angular testado sob o mesmo modelo.

\textbf{Perguntas:} a inclinação tem qual unidade? O intercepto em 0 m é ecologicamente interpretável? A topografia causa a diversidade? Uma nova localidade estará na faixa azul ou na faixa cinza?

\textbf{Resposta esperada:} Shannon por metro de DP da elevação. O zero fica fora do intervalo observado (4,4 a 30 m), limitando a interpretação do intercepto. O estudo é observacional. A faixa azul é IC da média; a cinza é IP de uma nova observação, mais amplo por incorporar o erro residual individual.

\textbf{Comentário de diagnóstico:} o local L04 tem a maior distância de Cook no modelo simples. Investigue o contexto e a estrutura entre ambientes. Não exclua o local para melhorar a significância. Uma observação influente pode revelar uma variável explicativa omitida, não um erro de campo.

\section{Atividade 7 --- Separar associação simples de associação parcial}

\textbf{Execução:} seções 9 e 10; figuras 14--16.

O modelo usa topografia centrada em 15 m e expressa por 10 m; chuva centrada em 2400 mm e expressa por 100 mm; Conservado é a referência do fator ambiente. O intercepto deixa de se referir a topografia e chuva iguais a zero, mas continua condicionado à categoria de referência e ao suporte dos dados.

\textbf{Hipótese parcial:} \ensuremath{\beta}topografia = 0 depois de incluir chuva e ambiente. Compare o modelo sem topografia com o modelo completo sobre as mesmas 36 linhas.

\textbf{Resultados:} coeficiente de 0,26867 Shannon por 10 m, IC95\% 0,18384 a 0,35350; F parcial(1,31) \ensuremath{\approx} 41,721, p \ensuremath{\approx} 3,35 \ensuremath{\times} 10\textsuperscript{-}\textsuperscript{7}. Ambiente é testado conjuntamente com dois gl: F(2,31) \ensuremath{\approx} 21,363, p \ensuremath{\approx} 1,47 \ensuremath{\times} 10\textsuperscript{-}\textsuperscript{6}.

\textbf{Ênfase:} \texttt{anova(um\_modelo)} produz uma decomposição sequencial. A ordem dos termos pode importar. \texttt{anova(reduzido, completo)} explicita a pergunta parcial, desde que os modelos sejam aninhados, sobre as mesmas observações e com a estrutura residual apropriada. O teste do fator inteiro não é o teste de apenas um dos seus coeficientes.

\textbf{Leitura do gráfico condicional:} linhas são médias previstas com chuva fixada em 2400 mm; os pontos tiveram apenas o termo de chuva removido. As linhas ficam dentro da faixa de topografia observada em cada ambiente. As inclinações são iguais por construção do modelo aditivo; não são três inclinações estimadas separadamente.

\textbf{Extensão de interação:} adicionar topografia \ensuremath{\times} ambiente produz F \ensuremath{\approx} 1,770 e p \ensuremath{\approx} 0,188. Isso não demonstra identidade das inclinações. Trata-se de uma extensão exploratória identificada, e não de justificativa para selecionar automaticamente a história ecológica pelo p.

\textbf{HC3:} compare os IC robustos e clássicos. Os coeficientes não mudam, mas a estimativa da covariância muda. HC3 com referência normal é assintótico e não equivale a corrigir dependência espacial. Para topografia, o IC robusto é aproximadamente 0,176 a 0,362 por 10 m.

\section{Atividade 8 --- Explicar o que “modelo melhor” significa}

\textbf{Execução:} seção 11; figura 17.

\begin{table}[H]\centering\small
\begin{tabular}{p{0.27\linewidth}p{0.23\linewidth}p{0.18\linewidth}p{0.2\linewidth}}
\toprule
Modelo & R\textsuperscript{2} ajustado & RMSE de ajuste & RMSE LOO\\
\midrule
Simples & 0,353 & 0,237 & 0,249\\
Múltiplo & 0,714 & 0,150 & 0,175\\
\bottomrule\end{tabular}\end{table}

O erro de ajuste usa as observações que estimaram os coeficientes; LOO omite cada local antes de prever sua resposta. Ainda é uma avaliação interna ao conjunto simulado, não validação regional independente. Para dados espaciais ou temporais dependentes, a separação de treino e avaliação precisa refletir essa dependência.

A inclinação parcial, ao omitir um local por vez, varia aproximadamente de 0,246 a 0,293 por 10 m. A figura mostra sensibilidade, não um ranking de locais a excluir. Discuta também se as métricas respondem à pergunta principal: explicação de associação e previsão são objetivos relacionados, mas diferentes.

\section{Extensão experimental --- DBC}

\textbf{Execução:} seção 12; figuras 18 e 19.

Há seis blocos completos e uma unidade por tratamento em cada bloco. F de tratamento ajustado para bloco \ensuremath{\approx} 105,428, com 2 e 10 gl. As linhas do perfil conectam unidades diferentes do mesmo bloco; não representam medidas repetidas no tempo nem uma dose quantitativa contínua.

Sem réplicas dentro da célula bloco \ensuremath{\times} tratamento, a interação não é separada do erro. O DBC e o DIC têm bases de experimentos diferentes: não compare seus quadrados médios como demonstração controlada da eficiência do bloqueamento.

\section{Extensão conceitual --- randomização}

\textbf{Execução:} seção 5; figura 07.

Das 4999 realocações, nenhuma atingiu o F observado. O p Monte Carlo corrigido é 1/5000 = 0,0002; não é zero e não resolve uma cauda da ordem de 10\textsuperscript{-}\textsuperscript{7}. A nula aguda afirma ausência de efeito em cada unidade; não é apenas igualdade das médias. A permutação livre de rótulos é justificada aqui pelo DIC completamente casualizado, não pelo fato de o comando \texttt{sample()} estar disponível.

\section{Fechamento --- produto da turma}

Peça uma síntese de até oito linhas que inclua: desenho e unidade; hipótese; diferença Alta - Baixa com IC simultâneo; associação parcial de topografia com IC; uma observação de diagnóstico; e uma limitação de inferência. Não aceite uma conclusão formada somente por valores-p.

O arquivo \texttt{saidas/comentarios\_resultados.txt} contém uma síntese calculada pelo script. Use-o como referência para discutir uma boa redação; não como texto a copiar para qualquer dissertação. Dados simulados demonstram o procedimento, e não resultados de campo.
